\documentclass[11pt]{article}

\usepackage[margin=1in]{geometry}
\usepackage{amsmath, amssymb, amsthm}
\usepackage{booktabs}
\usepackage{array}
\usepackage{enumitem}
\usepackage[colorlinks=true, linkcolor=blue, urlcolor=blue]{hyperref}
\usepackage{longtable}

\newcommand{\R}{\mathbb{R}}
\newcommand{\B}{\{0,1\}}
\DeclareMathOperator*{\maximize}{maximize}
\DeclareMathOperator*{\minimize}{minimize}

\theoremstyle{remark}
\newtheorem{remark}{Remark}[section]

\title{Person-to-Project Assignment as a Mixed-Integer Linear Program:\\
A Progressive Family of Formulations}
\author{}
\date{\today}

\begin{document}
\maketitle

\begin{abstract}
We formulate the problem of assigning people to projects as a mixed-integer
linear program (MILP). Rather than a single monolithic model, we develop a
\emph{progression of variants}: a static binary assignment model with hard
skill-coverage constraints (Variant~1); a time- and role-indexed model with
continuous hour allocations (Variant~2); and a model with soft coverage,
project selection, and shortfall penalties (Variant~3). We then present
cross-cutting constraint modules --- pairwise compatibility, team composition
and diversity, fairness, development goals, location --- that can be attached
to any variant, and close with a discussion of multi-objective design and
computational considerations.
\end{abstract}

\tableofcontents

%=====================================================================
\section{Problem Statement}
%=====================================================================

An organization must staff a portfolio of projects from a pool of people.
Each project demands certain skills at minimum proficiency levels, a
headcount per role, and effort over a planning horizon; it carries a
strategic value, a budget, and a timeline. Each person has skill
proficiencies, limited capacity, a cost rate, a seniority level,
preferences over projects, and a bound on concurrent engagements.
The planner must decide \emph{who works on what} (and, in richer variants,
\emph{in which role, in which period, for how many hours}) so that every
project's requirements are \emph{covered}, while trading off skill fit,
preference satisfaction, cost, and strategic value.

The defining structural feature is the \textbf{coverage constraint}: a
project's required skill set must be covered by the \emph{union} of the
skills of its assigned team, not by any single person. This makes the
problem a generalization of set covering intertwined with a generalized
assignment problem, and it is NP-hard in general.

%=====================================================================
\section{Notation}
%=====================================================================

\subsection{Sets and indices}

\begin{center}
\begin{tabular}{ll}
\toprule
Symbol & Meaning \\
\midrule
$i \in I$ & people \\
$j \in J$ & projects \\
$s \in S$ & skills \\
$r \in R$ & roles (e.g.\ lead, contributor, reviewer) \\
$t \in T$ & time periods (weeks, sprints, months) \\
$d \in D$ & teams / departments; $I_d \subseteq I$ is the set of people in $d$ \\
$l \in L$ & seniority levels (ordered) \\
$J^{\mathrm{man}} \subseteq J$ & mandatory (must-staff) projects \\
$T_j \subseteq T$ & active periods of project $j$ (from $\mathit{start}_j$ to $\mathit{end}_j$) \\
\bottomrule
\end{tabular}
\end{center}

\subsection{Person parameters}

\begin{center}
\begin{tabular}{ll}
\toprule
Symbol & Meaning \\
\midrule
$\mathit{prof}_{is} \in \{0,\dots,5\}$ & proficiency of person $i$ in skill $s$ \\
$\mathit{cap}_{it} \ge 0$ & available hours (or FTE fraction) of $i$ in period $t$ \\
$c_i \ge 0$ & cost rate of $i$ (per hour) \\
$\mathit{lev}_i \in L$ & seniority level of $i$ \\
$\mathit{tz}_i$ & time zone / location of $i$ \\
$p_{ij} \in \R$ & preference (interest) score of $i$ for project $j$ \\
$\mathit{cert}_{ik} \in \B$ & person $i$ holds certification/clearance $k$ \\
$\overline{n}_i \in \mathbb{Z}_{+}$ & max number of concurrent projects for $i$ \\
$g_{is} \in \B$ & $i$ wants to \emph{grow} skill $s$ (development goal) \\
\bottomrule
\end{tabular}
\end{center}

Existing commitments are folded into $\mathit{cap}_{it}$ (capacity net of
prior obligations). Language proficiency and domain/client experience are
treated as additional ``skills'' in $S$, which keeps the model uniform.

\subsection{Project parameters}

\begin{center}
\begin{tabular}{ll}
\toprule
Symbol & Meaning \\
\midrule
$\pi_{js} \in \{0,\dots,5\}$ & minimum proficiency in skill $s$ required by $j$ ($0$ = not needed) \\
$n_{js} \in \mathbb{Z}_{+}$ & number of assigned people who must meet skill $s$ (coverage depth) \\
$\delta_{jr} \in \mathbb{Z}_{+}$ & required headcount in role $r$ on project $j$ \\
$e_{jt} \ge 0$ & effort (hours) needed by $j$ in period $t$ \\
$v_j \ge 0$ & strategic value / priority of $j$ \\
$B_j \ge 0$ & budget cap of $j$ \\
$\underline{N}_j, \overline{N}_j$ & min / max team size of $j$ \\
$\lambda_j \in L$ & minimum seniority required for the lead role on $j$ \\
$\mathit{creq}_{jk} \in \B$ & project $j$ requires certification/clearance $k$ \\
\bottomrule
\end{tabular}
\end{center}

\subsection{Relational parameters}

\begin{center}
\begin{tabular}{ll}
\toprule
Symbol & Meaning \\
\midrule
$f_{ij} \in \R$ & skill-fit score of person $i$ for project $j$ (precomputed) \\
$P^{-} \subseteq I \times I$ & pairs who \emph{cannot} work together \\
$P^{+} \subseteq I \times I$ & pairs who \emph{must} work together (if either is assigned) \\
$\sigma_{i i'} \in \R$ & synergy score for pairing $i$ and $i'$ (e.g.\ prior collaboration) \\
$\rho_d \in (0,1]$ & max fraction of any team drawn from department $d$ \\
\bottomrule
\end{tabular}
\end{center}

\begin{remark}[Derived eligibility sets]
\label{rem:eligibility}
Much of the combinatorics is best handled by \emph{preprocessing}. Define
\begin{align*}
Q_{js} &= \{\, i \in I : \mathit{prof}_{is} \ge \pi_{js} \,\}
  && \text{(people qualified for skill $s$ on project $j$),} \\
E_j &= \{\, i \in I : \mathit{cert}_{ik} \ge \mathit{creq}_{jk}\ \forall k,
      \ \text{location/tz compatible with } j \,\}
  && \text{(people eligible for project $j$).}
\end{align*}
Variables are only created for $(i,j)$ with $i \in E_j$; ineligible
assignments are never instantiated. This is both a correctness device and
the single most effective way to keep instances small.
\end{remark}

%=====================================================================
\section{Variant 1: Static Assignment with Hard Coverage}
%=====================================================================

The base model answers one question: \emph{which people staff which
projects?} Time, roles, and hours are aggregated away; each project $j$
exposes an estimated per-person workload $w_j$ (hours over its lifetime),
so the cost of assigning $i$ to $j$ is $c_{ij} = c_i \, w_j$.

\subsection{Decision variables}
\begin{align*}
x_{ij} &\in \B && \text{person $i$ is assigned to project $j$}
  \quad (i \in E_j).
\end{align*}

\subsection{Formulation}

\begin{align}
\maximize_{x} \quad
  & \alpha \sum_{j \in J} \sum_{i \in E_j} f_{ij}\, x_{ij}
  \;+\; \beta \sum_{j \in J} \sum_{i \in E_j} p_{ij}\, x_{ij}
  \;-\; \gamma \sum_{j \in J} \sum_{i \in E_j} c_{ij}\, x_{ij}
  \label{eq:v1-obj} \\[4pt]
\text{s.t.} \quad
  & \sum_{i \in Q_{js} \cap E_j} x_{ij} \;\ge\; n_{js}
  && \forall j \in J,\ s \in S:\ \pi_{js} \ge 1
  \label{eq:v1-cover} \\
  & \underline{N}_j \;\le\; \sum_{i \in E_j} x_{ij} \;\le\; \overline{N}_j
  && \forall j \in J
  \label{eq:v1-teamsize} \\
  & \sum_{j:\, i \in E_j} x_{ij} \;\le\; \overline{n}_i
  && \forall i \in I
  \label{eq:v1-concurrency} \\
  & \sum_{i \in E_j} c_{ij}\, x_{ij} \;\le\; B_j
  && \forall j \in J
  \label{eq:v1-budget} \\
  & x_{ij} \in \B
  && \forall j \in J,\ i \in E_j. \nonumber
\end{align}

Constraint \eqref{eq:v1-cover} is the \textbf{skill-coverage} constraint:
for every skill $s$ that project $j$ requires, at least $n_{js}$ of the
assigned people must be qualified in $s$ (the common case is
$n_{js} = 1$; setting $n_{js} = 2$ buys redundancy for critical skills).
Constraint \eqref{eq:v1-teamsize} bounds team size,
\eqref{eq:v1-concurrency} caps each person's concurrent engagements, and
\eqref{eq:v1-budget} enforces project budgets.

\subsection{Explicit coverage variables (skill-load control)}
\label{sec:v1-explicit}

Formulation \eqref{eq:v1-cover} lets a single versatile person
``count'' toward every skill simultaneously, which may be unrealistic: one
person cannot in practice \emph{own} five skill areas on one project. To
cap the number of skill requirements a person may cover, introduce
\begin{align*}
y_{ijs} \in \B \quad \text{: person $i$ is designated to cover skill $s$
on project $j$} \quad (i \in Q_{js} \cap E_j),
\end{align*}
and replace \eqref{eq:v1-cover} with
\begin{align}
\sum_{i \in Q_{js} \cap E_j} y_{ijs} &\;\ge\; n_{js}
  && \forall j,\ s:\ \pi_{js} \ge 1
  \label{eq:v1b-cover} \\
y_{ijs} &\;\le\; x_{ij}
  && \forall j,\ s,\ i \in Q_{js} \cap E_j
  \label{eq:v1b-link} \\
\sum_{s:\, i \in Q_{js}} y_{ijs} &\;\le\; \kappa
  && \forall j,\ i \in E_j,
  \label{eq:v1b-load}
\end{align}
where $\kappa$ is the maximum number of skill areas one person may own on a
single project. This pattern --- an explicit ``who covers what'' layer under
the assignment layer --- recurs whenever coverage must be attributed rather
than merely counted.

\begin{remark}[When Variant 1 suffices]
Use Variant~1 for annual or quarterly portfolio staffing where projects are
long-lived relative to the planning cycle, workloads are roughly uniform,
and the interesting decision is team \emph{composition}, not scheduling.
\end{remark}

%=====================================================================
\section{Variant 2: Time- and Role-Indexed Allocation}
%=====================================================================

When projects have effort profiles over time, people have fluctuating
availability, and roles matter (a project needs \emph{a lead} and
\emph{two contributors}, not just three bodies), the assignment must be
indexed by role and period, and hours become a continuous decision.

\subsection{Decision variables}
\begin{align*}
x_{ijrt} &\in \B
  && \text{person $i$ works on project $j$ in role $r$ during period $t$}
  \quad (i \in E_j,\ t \in T_j), \\
h_{ijt} &\ge 0
  && \text{hours of person $i$ allocated to project $j$ in period $t$}, \\
a_{ijt} &\in \B
  && \text{person $i$ is active on project $j$ in period $t$
      (aggregation of roles)}, \\
z_{ij} &\in \B
  && \text{person $i$ is \emph{ever} assigned to project $j$
      (membership indicator)}.
\end{align*}
The auxiliary variables $a_{ijt}$ and $z_{ij}$ are defined by linear
linking constraints below; they let concurrency, coverage, and continuity
be stated cleanly.

\subsection{Formulation}

\begin{align}
\maximize \quad
  & \alpha \sum_{i,j} f_{ij}\, z_{ij}
  \;+\; \beta \sum_{i,j} p_{ij}\, z_{ij}
  \;-\; \gamma \sum_{i,j,t} c_i\, h_{ijt}
  \;-\; \mu \sum_{i,j} z_{ij}
  \label{eq:v2-obj} \\[4pt]
\text{s.t.}\quad
% role structure
  & \sum_{r \in R} x_{ijrt} \;=\; a_{ijt}
  && \forall i,\ j,\ t \in T_j
  \label{eq:v2-onerole} \\
  & a_{ijt} \;\le\; z_{ij}
  && \forall i,\ j,\ t \in T_j
  \label{eq:v2-membership} \\
% role coverage
  & \sum_{i \in E_j} x_{ijrt} \;\ge\; \delta_{jr}
  && \forall j,\ r,\ t \in T_j
  \label{eq:v2-rolecover} \\
% skill coverage each period
  & \sum_{i \in Q_{js} \cap E_j} a_{ijt} \;\ge\; n_{js}
  && \forall j,\ s:\ \pi_{js}\ge 1,\ t \in T_j
  \label{eq:v2-skillcover} \\
% seniority mix
  & \sum_{i:\, \mathit{lev}_i \ge \lambda_j} x_{ij,\mathrm{lead},t} \;\ge\; 1
  && \forall j,\ t \in T_j
  \label{eq:v2-seniority} \\
% hours vs assignment
  & h_{ijt} \;\le\; \mathit{cap}_{it}\; a_{ijt}
  && \forall i,\ j,\ t \in T_j
  \label{eq:v2-hours-ub} \\
  & h_{ijt} \;\ge\; \underline{h}\; a_{ijt}
  && \forall i,\ j,\ t \in T_j
  \label{eq:v2-hours-lb} \\
% capacity
  & \sum_{j:\, t \in T_j} h_{ijt} \;\le\; \mathit{cap}_{it}
  && \forall i,\ t
  \label{eq:v2-capacity} \\
% effort demand
  & \sum_{i \in E_j} h_{ijt} \;\ge\; e_{jt}
  && \forall j,\ t \in T_j
  \label{eq:v2-effort} \\
% concurrency per period
  & \sum_{j:\, t \in T_j} a_{ijt} \;\le\; \overline{n}_i
  && \forall i,\ t
  \label{eq:v2-concurrency} \\
% budget over lifetime
  & \sum_{t \in T_j} \sum_{i \in E_j} c_i\, h_{ijt} \;\le\; B_j
  && \forall j
  \label{eq:v2-budget} \\
  & x_{ijrt},\, a_{ijt},\, z_{ij} \in \B, \quad h_{ijt} \ge 0. \nonumber
\end{align}

\paragraph{Reading the model.}
\eqref{eq:v2-onerole} says a person active on a project in a period holds
exactly one role there. \eqref{eq:v2-rolecover} is per-period role coverage
($\delta_{jr}$ can be made time-dependent $\delta_{jrt}$ if a project needs
a reviewer only near delivery). \eqref{eq:v2-skillcover} re-imposes skill
coverage \emph{in every active period}, so coverage cannot lapse when a
team rotates. \eqref{eq:v2-seniority} guarantees the lead in each period is
sufficiently senior. \eqref{eq:v2-hours-ub} links hours to assignment with
the \emph{tight} big-M $\mathit{cap}_{it}$ (never use a generic large $M$
here); \eqref{eq:v2-hours-lb} imposes a minimum meaningful allocation
$\underline{h}$ (e.g.\ 4 hours/week) to forbid sliver assignments.
The term $-\mu \sum z_{ij}$ in \eqref{eq:v2-obj} charges a small fixed cost
per team membership, discouraging fragmented teams; $\mu$ also controls
the classic trade-off between a few heavily-loaded generalists and many
lightly-loaded specialists.

\subsection{Continuity and churn control}
\label{sec:continuity}

Re-solving a time-indexed model can churn assignments. Two standard,
composable devices:
\begin{enumerate}[label=(\alph*)]
\item \textbf{Team stability within the horizon.} Penalize onboarding events with
  $u_{ijt} \in \B$ ($u \ge a_{ijt} - a_{ij,t-1}$) and add
  $-\theta \sum u_{ijt}$ to the objective: each time a person \emph{joins} a
  project mid-stream costs $\theta$.
\item \textbf{Stability across re-solves.} Given the incumbent plan
  $\hat{a}_{ijt}$, penalize deviations
  $\sum_{i,j,t} \tau_{ijt} \, | a_{ijt} - \hat{a}_{ijt} |$, linearized in the
  usual way; or hard-fix assignments inside a frozen window
  $t \le t_{\mathrm{freeze}}$.
\end{enumerate}

\begin{remark}[When Variant 2 is worth its size]
The variable count is $O(|I|\,|J|\,|R|\,|T|)$. Use it when effort profiles
$e_{jt}$ genuinely vary, when part-time multi-project allocation is the norm
(matrix organizations, consultancies), or when role hand-offs over time are
part of the decision. If everyone works one project at a time full-time,
Variant~1 plus a start-date heuristic is usually better value.
\end{remark}

%=====================================================================
\section{Variant 3: Soft Coverage, Project Selection, and Penalties}
%=====================================================================

Variants 1--2 presume every project must be staffed and can be. In
practice demand exceeds supply: the model should \emph{choose} which
optional projects to run, fully staff the mandatory ones, and degrade
gracefully --- reporting \emph{where} coverage falls short rather than
returning \textsc{infeasible}.

\subsection{Additional decision variables}
\begin{align*}
w_j &\in \B && \text{project $j$ is selected (staffed at all)}, \\
u_{js} &\ge 0 && \text{shortfall in coverage of skill $s$ on project $j$}, \\
q_{jt} &\ge 0 && \text{shortfall in effort on project $j$ in period $t$
  \ (time-indexed setting)}.
\end{align*}

\subsection{Formulation (built on Variant 1; the same surgery applies to Variant 2)}

\begin{align}
\maximize \quad
  & \sum_{j} v_j\, w_j
  \;+\; \alpha \sum_{i,j} f_{ij}\, x_{ij}
  \;+\; \beta \sum_{i,j} p_{ij}\, x_{ij}
  \;-\; \gamma \sum_{i,j} c_{ij}\, x_{ij}
  \;-\; \sum_{j,s} \phi_{js}\, u_{js}
  \label{eq:v3-obj} \\[4pt]
\text{s.t.}\quad
  & \sum_{i \in Q_{js} \cap E_j} x_{ij} \;\ge\; n_{js}\, w_j \;-\; u_{js}
  && \forall j,\ s:\ \pi_{js} \ge 1
  \label{eq:v3-softcover} \\
  & u_{js} \;\le\; \overline{u}_{js}
  && \forall j,\ s
  \label{eq:v3-slackcap} \\
  & x_{ij} \;\le\; w_j
  && \forall j,\ i \in E_j
  \label{eq:v3-link} \\
  & \underline{N}_j\, w_j \;\le\; \sum_{i \in E_j} x_{ij}
    \;\le\; \overline{N}_j\, w_j
  && \forall j
  \label{eq:v3-teamsize} \\
  & w_j \;=\; 1
  && \forall j \in J^{\mathrm{man}}
  \label{eq:v3-mandatory} \\
  & \text{\eqref{eq:v1-concurrency}, \eqref{eq:v1-budget}},
  \quad x_{ij},\, w_j \in \B,\quad u_{js} \ge 0. \nonumber
\end{align}

\paragraph{Design notes.}
\begin{itemize}
\item \textbf{Selection.} An unselected project ($w_j = 0$) absorbs no
  people, by \eqref{eq:v3-link} and \eqref{eq:v3-teamsize}, and forfeits its
  value $v_j$. The model thus performs portfolio triage and staffing
  \emph{jointly}, which dominates the common two-phase practice of first
  picking projects and then discovering they cannot be staffed.
\item \textbf{Soft coverage.} The slack $u_{js}$ converts a hard coverage
  row into a priced one. Penalties should be \emph{tiered}:
  $\phi_{js} \gg \max_j v_j$ for mission-critical skills (effectively hard),
  moderate for nice-to-have depth. The cap \eqref{eq:v3-slackcap} prevents
  the absurd solution of ``staffing'' a selected project entirely with slack
  --- e.g.\ $\overline{u}_{js} = n_{js} - 1$ forces at least one qualified
  person on every required skill of a selected project.
\item \textbf{Mandatory projects} keep hard coverage simply by fixing
  $\overline{u}_{js} = 0$ for $j \in J^{\mathrm{man}}$.
\item \textbf{Diagnosis.} At optimum, the positive $u_{js}$ are exactly the
  hiring/training gaps: read them as a ranked shopping list for
  recruitment, with shadow value $\phi_{js}$ per unit.
\end{itemize}

In the time-indexed setting the same treatment applies to effort:
replace \eqref{eq:v2-effort} by
$\sum_i h_{ijt} \ge e_{jt}\, w_j - q_{jt}$ with penalty
$\sum_{j,t} \psi_j\, q_{jt}$, capturing schedule slippage rather than
skill gaps.

%=====================================================================
\section{Cross-Cutting Constraint Modules}
%=====================================================================

The following modules attach to any variant. Throughout, read $x_{ij}$ as
$z_{ij}$ in the time-indexed setting.

\subsection{Pairwise compatibility}
\begin{align}
x_{ij} + x_{i'j} &\;\le\; 1
  && \forall (i,i') \in P^{-},\ \forall j
  \label{eq:cannot} \\
x_{ij} &\;=\; x_{i'j}
  && \forall (i,i') \in P^{+},\ \forall j.
  \label{eq:must}
\end{align}
Constraint \eqref{eq:cannot} covers interpersonal conflicts and
manager--report separation (instantiate $P^{-}$ from the org chart when
policy forbids a manager reviewing their own report's project work). The
must-pair constraint \eqref{eq:must} is strong (both or neither); the
weaker ``if $i$ then $i'$'' is $x_{ij} \le x_{i'j}$, useful for
apprentice--mentor pairings.

\subsection{Synergy bonuses (linearized products)}
To \emph{reward} (not require) putting a productive pair together, add
$w_{i i' j} \in [0,1]$ with the McCormick equalities for a product of
binaries:
\begin{align}
w_{i i' j} \le x_{ij}, \qquad
w_{i i' j} \le x_{i'j}, \qquad
w_{i i' j} \ge x_{ij} + x_{i'j} - 1,
\label{eq:mccormick}
\end{align}
and the objective term $+\sum \sigma_{i i'}\, w_{i i' j}$. With
$\sigma_{i i'} > 0$ the first two inequalities suffice (the maximization
pushes $w$ up); include the third only for negative synergies. Instantiate
these variables \emph{sparsely} --- only for pairs with
$|\sigma_{i i'}|$ above a threshold --- or the model inflates
quadratically.

\subsection{Team composition and diversity}
A ratio cap ``no more than a fraction $\rho_d$ of any team from department
$d$'' looks nonlinear (team size is a variable) but linearizes exactly:
\begin{align}
\sum_{i \in I_d \cap E_j} x_{ij}
  \;\le\; \rho_d \sum_{i \in E_j} x_{ij}
  && \forall j,\ d.
  \label{eq:diversity}
\end{align}
Floors ($\ge$) work symmetrically, e.g.\ representation targets across
departments or sites. A minimum seniority \emph{mix} beyond the lead rule
is the same pattern with $I_d$ replaced by $\{i : \mathit{lev}_i \ge l\}$.

\subsection{Location and time-zone overlap}
Discretize into zone clusters $Z$, let $I_z$ be people in cluster $z$, and
require a quorum in a core cluster $z^{*}_j$ chosen per project:
\begin{align}
\sum_{i \in I_{z^{*}_j} \cap E_j} x_{ij}
  \;\ge\; \eta_j \sum_{i \in E_j} x_{ij},
\end{align}
i.e.\ at least a fraction $\eta_j$ of the team shares core working hours.
Hard on-site requirements are handled upstream via the eligibility sets
$E_j$ (Remark~\ref{rem:eligibility}), not as constraints.

\subsection{Fairness of load}
Unconstrained cost minimization concentrates work on the cheapest capable
people. Two remedies:
\begin{enumerate}[label=(\alph*)]
\item \textbf{Min--max utilization.} With utilization
  $\mathit{util}_i = \sum_{j,t} h_{ijt} \big/ \sum_t \mathit{cap}_{it}$,
  add $\Lambda \ge \mathit{util}_i \ \forall i$ and the term
  $-\omega \Lambda$ to the objective.
\item \textbf{Balanced project counts.} Bound the spread:
  $\underline{\nu} \le \sum_j x_{ij} \le \overline{\nu}$ for all $i$ in a
  designated pool (everyone gets between $\underline{\nu}$ and
  $\overline{\nu}$ assignments).
\end{enumerate}

\subsection{Development goals (stretch assignments)}
Growth is exposure to a desired skill \emph{without} being counted on for
coverage. Let $G_{js} \subseteq I$ be people with $g_{is} = 1$ who are
\emph{near}-qualified ($\mathit{prof}_{is} = \pi_{js} - 1$). Introduce
$y^{g}_{ijs} \in \B$ (a stretch slot) with
\begin{align}
y^{g}_{ijs} \;\le\; x_{ij}, \qquad
\sum_{s} y^{g}_{ijs} \;\le\; 1, \qquad
\sum_{i \in G_{js}} y^{g}_{ijs} \;\le\; 1
  \quad \forall j, s,
\end{align}
and the objective bonus $+\epsilon \sum y^{g}_{ijs}$. Crucially,
$y^{g}$ does \emph{not} appear in the coverage constraints: the stretch
person shadows a qualified one (who is still required by
\eqref{eq:v1-cover}). If desired, couple them:
$y^{g}_{ijs} \le \sum_{i' \in Q_{js}} y_{i'js}$ using the explicit coverage
variables of Section~\ref{sec:v1-explicit}, so a stretch slot exists only
where a qualified owner does.

%=====================================================================
\section{Objective Design and Multi-Objective Trade-offs}
%=====================================================================

The criteria in play --- skill fit, preference satisfaction, cost,
strategic value covered, plus soft penalties --- are incommensurable. The
formulations above use a \textbf{weighted sum}
\begin{align}
\maximize\quad
  \alpha\, \mathrm{Fit} + \beta\, \mathrm{Pref}
  + \sum_j v_j w_j - \gamma\, \mathrm{Cost}
  - \textstyle\sum \phi\, u - \textstyle\sum \psi\, q - \dots
\label{eq:weighted}
\end{align}
which is the practical default, but the weights deserve care.

\subsection{Normalization}
Raw scales differ wildly (fit scores in $[0,1]$, costs in the
$10^{5}$--$10^{6}$ range), so raw weights are meaningless. Normalize each
term by its \emph{ideal value}: solve the single-objective problem
$\max \mathrm{Fit}$ to get $\mathrm{Fit}^{*}$, likewise
$\mathrm{Cost}^{*} = \min \mathrm{Cost}$ over feasible staffings, then
optimize
\begin{align}
\alpha\, \frac{\mathrm{Fit}}{\mathrm{Fit}^{*}}
- \gamma\, \frac{\mathrm{Cost}}{\mathrm{Cost}^{*}} + \dots
\end{align}
so that $\alpha, \gamma$ express genuine relative importance. Each ideal
value costs one auxiliary MILP solve --- cheap, and reusable across
weight settings.

\subsection{Lexicographic (preemptive) ordering}
When priorities are strict --- e.g.\ (1) staff all mandatory projects with
zero coverage shortfall, (2) maximize strategic value of optional projects,
(3) minimize cost, (4) maximize preferences --- solve a sequence of MILPs:
optimize objective 1; add the constraint
$\text{obj}_1 \ge (1 - \varepsilon)\, \text{obj}_1^{*}$; optimize
objective 2 over the restricted set; and so on. The tolerance
$\varepsilon$ (e.g.\ 1\%) trades a sliver of a higher priority for
substantial freedom in lower ones and markedly better solutions in
practice than $\varepsilon = 0$.

\subsection{$\varepsilon$-constraint method}
Keep one objective, bound the rest:
\begin{align}
\maximize \ \mathrm{Fit} \quad \text{s.t.} \quad
\mathrm{Cost} \le \overline{C}, \qquad
\sum_j v_j w_j \ge \underline{V},
\end{align}
then sweep $\overline{C}$ (and/or $\underline{V}$) to trace the Pareto
frontier. This is the method of choice for \emph{presenting} the trade-off
to decision-makers: ``every additional \$50k of staffing budget buys this
much fit'' is a curve, not a weight.

\subsection{Goal programming}
When stakeholders state targets (``average utilization near 80\%,''
``preference score at least 0.7 per person''), attach deviation variables
to each goal, $\mathrm{goal}_k - d^{+}_k + d^{-}_k = \mathrm{target}_k$,
and minimize the weighted sum of the undesirable deviations. This
subsumes the soft-coverage device of Variant~3, and it is the honest
formulation when management thinks in targets rather than trade-off rates.

\subsection{Which to use}
Weighted sum for routine automated runs (fast, single solve); lexicographic
when governance imposes a strict priority order; $\varepsilon$-constraint
for stakeholder-facing what-if analysis; goal programming when the input
arrives as targets. They can be mixed: lexicographic at the top (mandatory
coverage first), weighted sum within the final level.

%=====================================================================
\section{Computational Considerations}
%=====================================================================

\begin{itemize}
\item \textbf{Scale.} Variant 1 has $O(|I|\,|J|)$ binaries --- typically
  thousands, easy for modern solvers. Variant 2 has
  $O(|I|\,|J|\,|R|\,|T|)$; at $200$ people $\times$ $40$ projects $\times$
  $3$ roles $\times$ $26$ weeks this is $\sim\!6 \times 10^{5}$ binaries,
  demanding but tractable \emph{if} eligibility preprocessing
  (Remark~\ref{rem:eligibility}) prunes the majority of $(i,j)$ pairs, as
  it usually does.
\item \textbf{Tight big-Ms.} The only big-M in the family is
  \eqref{eq:v2-hours-ub}, and $\mathit{cap}_{it}$ is its tightest possible
  value. Resist the generic $M = 10^{6}$; the LP relaxation gap it opens is
  the difference between minutes and hours of solve time.
\item \textbf{Symmetry.} Interchangeable people (identical skills, cost,
  availability) create symmetric solutions that bloat the branch-and-bound
  tree. Either aggregate them into pools with integer ``count'' variables,
  or add lexicographic ordering cuts $x_{ij} \ge x_{i'j}$-style over each
  symmetry class. Role \emph{slots} should never be modeled individually
  (``contributor slot 1, slot 2'') --- the headcount form
  \eqref{eq:v2-rolecover} is symmetry-free by construction.
\item \textbf{Rolling horizon.} For long horizons, solve Variant 2 over a
  window (e.g.\ 8 weeks detailed + remainder aggregated to Variant-1
  granularity), commit the first weeks, roll forward. Combine with the
  re-solve stability penalty of Section~\ref{sec:continuity}.
\item \textbf{Warm starts.} The incumbent staffing plan is almost always
  feasible for the re-solve after minor data changes; pass it as a MIP
  start.
\item \textbf{Infeasibility handling.} Prefer Variant 3's priced slacks to
  post-hoc IIS analysis: the optimal slack pattern is itself the
  diagnosis, in the units managers care about.
\end{itemize}

%=====================================================================
\section{Choosing a Variant: Summary}
%=====================================================================

\begin{center}
\begin{tabular}{p{3.4cm} p{4.6cm} p{6.2cm}}
\toprule
Variant & Decision granularity & Use when \\
\midrule
1. Static, hard coverage &
$x_{ij}$ &
Portfolio-level team composition; stable long-lived projects; staffing is
feasible and mandatory. \\[6pt]
2. Time- and role-indexed &
$x_{ijrt}$, $h_{ijt}$ &
Varying effort profiles; part-time multi-project work; roles and hand-offs
matter; capacity is the binding resource. \\[6pt]
3. Soft coverage + selection &
Variant 1 or 2 $+\ w_j,\ u_{js},\ q_{jt}$ &
Demand exceeds supply; portfolio triage is part of the question;
infeasibility must degrade into priced, diagnosable shortfalls. \\
\bottomrule
\end{tabular}
\end{center}

The modules of Section~5 --- pairwise compatibility, synergy, diversity
ratios, time-zone quorum, fairness, stretch assignments --- bolt onto any
row of this table, and the objective machinery of Section~6 applies
uniformly. In practice the recommended path is: start with Variant~1 plus
soft coverage (the Variant-3 surgery), validate the data and weights with
$\varepsilon$-constraint sweeps, and move to Variant~2 only for the
projects and people where period-level capacity is demonstrably the
binding constraint.

\end{document}
